% P03.1 rational collar minimax and boundary saturation proof, live integration.

\begin{proof}[\Cref{lem:p031-rational-protected-collar-gluing}]
The two displayed formulas agree on the overlap
$D\setminus\overline V$ by
\eqref{eq:p031-rational-protected-collar-order}, so they define a function
$H$ on the ordinary open cover
\[
                         X=D\cup(X\setminus\overline V).
\]
The $\mathbf Q$-LFS class is local by
\Cref{def:p031-rational-lfs}.  It is closed under finite maxima: locally
compress the two rational presentations as in
\eqref{eq:p031-rational-lfs-block-compression}, take a common denominator,
and combine the resulting finite families of analytic units.  Hence $H$ is
$\mathbf Q$-LFS.  Both $b$ and $u$ lie below $f$ on their domains, so the
same is true of $H$.
\end{proof}

\begin{proof}[\Cref{lem:p031-small-nonempty-boundary-domain}]
A separated good analytic space is locally compact Hausdorff and therefore
regular.  Since a connected Hausdorff space with at least two points has no
isolated point, the prescribed neighborhood contains a point $y\ne x$.
Regularity and local compactness give an ordinary open neighborhood $D$ of
$x$ whose compact closure is contained in the prescribed neighborhood
minus $\{y\}$.  If $\partial D$ were empty, then $D$ would be a nonempty
clopen proper subset of $X$, contradicting connectedness.  Thus
$\partial D\ne\varnothing$.
\end{proof}

\begin{proof}[\Cref{thm:p031-rational-collar-minimax}]
All functions in $\mathcal Q_f(U)$ are continuous.  A global
$\mathbf Q$-LFS minorant forces $f$ to be real-valued, and it also makes
$s_{\mathbf Q,X}$ and $P_{\mathbf Q}^{\mathrm{loc}}f$ finite.  In
particular, every maximum and difference used below is finite.

Fix an admissible $D$ and $b\in\mathcal Q_f(X)$.  We first prove
\begin{equation}
\label{eq:p031-rational-collar-lower-bound}
                       \max_{\partial D}(v-b)\geq\delta_x(v).
\end{equation}
Suppose instead that this maximum is smaller than $\delta_x(v)$.  Choose
$\gamma\in K^\times$ with $\tau:=-\log|\gamma|>0$.  Since
$\mathbf Q\tau$ is dense in $\mathbf R$, there is a number
\begin{equation}
\label{eq:p031-rational-collar-shift-choice}
 c=q\tau,\qquad q\in\mathbf Q_{>0},\qquad
 \max\bigl\{0,\max_{\partial D}(v-b)\bigr\}<c<\delta_x(v).
\end{equation}
The constant $-c$ is $\mathbf Q$-LFS because
\[
                   -c=q\max\{-\log|\gamma^{-1}|\}.
\]

Put $\phi:=v-b$ on $\overline D$.  Compactness of $\partial D$ and
continuity give a relative open neighborhood $N$ of $\partial D$ in
$\overline D$ on which $\phi<c$.  On the other hand,
\[
 \phi(x)=v(x)-b(x)\geq v(x)-s_{\mathbf Q,X}(x)
                         =\delta_x(v)>c,
\]
so $x\notin N$.  The compact set $\overline D\setminus N$ lies in $D$.
Choose an ordinary open $V$ with
\[
 x\in\overline D\setminus N\subset V,
                    \qquad\overline V\subset D.
\]
Then $v-c<b$ on $D\setminus\overline V$.  Apply
\Cref{lem:p031-rational-protected-collar-gluing} to $u:=v-c$.  It produces
a global $\mathbf Q$-LFS minorant $H$ of $f$ satisfying
\[
                    H(x)\geq v(x)-c>s_{\mathbf Q,X}(x),
\]
contrary to the definition of the global rational envelope.  This proves
\eqref{eq:p031-rational-collar-lower-bound} and hence
$\Lambda_x(v;W)\geq\delta_x(v)$.

For the reverse inequality, fix $\varepsilon>0$.  Choose
$b\in\mathcal Q_f(X)$ with
\[
                 b(x)>s_{\mathbf Q,X}(x)-\frac{\varepsilon}{3}.
\]
Continuity of $v-b$ gives an ordinary neighborhood $N$ of $x$, with
$N\subset W$, such that
\[
 (v-b)(z)<(v-b)(x)+\frac{\varepsilon}{3}
       <\delta_x(v)+\frac{2\varepsilon}{3}
                  \qquad(z\in N).
\]
By \Cref{lem:p031-small-nonempty-boundary-domain}, choose an admissible
$D$ with $\overline D\subset N$.  Then
\[
                  \max_{\partial D}(v-b)
                       <\delta_x(v)+\varepsilon.
\]
Taking the two infima and then letting $\varepsilon\downarrow0$ proves
\eqref{eq:p031-rational-collar-minimax}.

Finally, the definition of the local rational envelope gives
\[
 \Delta^{\mathbf Q}_{f,X}(x)
   =\sup_{\substack{x\in W\ \mathrm{ordinary\ open}\\
                     v\in\mathcal Q_f(W)}}
          \bigl(v(x)-s_{\mathbf Q,X}(x)\bigr).
\]
The left side is nonnegative.  If it is positive, the subfamily with
positive gap has the same supremum; on that subfamily
\eqref{eq:p031-rational-collar-minimax} replaces the gap by the collar
load.  If it is zero, the positive-gap subfamily is empty.  This proves
\eqref{eq:p031-rational-defect-collar-formula} with the stated convention.
\end{proof}

\begin{proof}[\Cref{cor:p031-rational-collar-absorption}]
Assume first that $\Delta^{\mathbf Q}_{f,X}(x)=0$.  Then every
$v\in\mathcal Q_f(W)$ satisfies $v(x)\leq s_{\mathbf Q,X}(x)$.  Given
$\varepsilon>0$, choose $b\in\mathcal Q_f(X)$ with
$b(x)>s_{\mathbf Q,X}(x)-\varepsilon/3$.  Continuity gives a neighborhood
$N\subset W$ of $x$ on which $v-b<2\varepsilon/3$.  Choose $D$ with
$\overline D\subset N$ and nonempty boundary by
\Cref{lem:p031-small-nonempty-boundary-domain}.  This gives
\eqref{eq:p031-rational-collar-absorption}.

Conversely, if the defect were positive, there would be a pair $(W,v)$
with $\delta_x(v)>0$.  By
\eqref{eq:p031-rational-collar-minimax}, every admissible $D$ and $b$ would
satisfy
\[
                         \max_{\partial D}(v-b)\geq\delta_x(v),
\]
contradicting \eqref{eq:p031-rational-collar-absorption} with
$\varepsilon:=\delta_x(v)/2$.
\end{proof}

\begin{proof}[\Cref{cor:p031-rational-boundary-saturation-collapse}]
For $v\in\mathcal Q_f(W)$ one has $v\leq f$, so rational boundary
saturation implies the collar-absorption condition
\eqref{eq:p031-rational-collar-absorption}.  Hence
\Cref{cor:p031-rational-collar-absorption} gives
$\Delta^{\mathbf Q}_{f,X}(x)=0$.  Apply the three-level ledger
\eqref{eq:p031-three-level-defect-ledger}.  It writes zero as sums of the
five nonnegative defects and gaps in
\eqref{eq:p031-rational-boundary-saturation-zero-ledger}; all of them
therefore vanish.  Their definitions and the local cofinality identities
\eqref{eq:p031-local-real-rational-lfs-cofinality} and
\Cref{thm:p031-local-pl-lfs-cofinality} give
\eqref{eq:p031-rational-boundary-saturation-collapse}.
\end{proof}

\begin{proof}[\Cref{cor:p031-rational-boundary-saturation-tower-collapse}]
By \Cref{cor:p031-rational-boundary-saturation-collapse},
$\Delta_0^{\mathbf Q}(x_0)=0$.  Given
$y\in\Pi_i^{-1}(x_0)$, its successive images form a compatible lift chain.
Apply \Cref{cor:p031-rational-cofinality-tower-collapse} to that chain.
This proves \eqref{eq:p031-rational-boundary-saturation-tower-zero} and
\eqref{eq:p031-rational-boundary-saturation-tower-collapse}.  Since $y$
was arbitrary, the conclusion holds on the complete inverse image at every
stage.
\end{proof}
